\section{Methods}
\label{sec:methods}

\subsection{Data}

We use two publicly available ASD GWAS summary-statistics releases:
\begin{enumerate}[nosep,leftmargin=1.2em]
  \item \textbf{iPSYCH--PGC ASD Nov.~2017} (Grove et al.~\cite{grove2019})
        with $N = 46\,350$ (18\,381 cases, 27\,969 controls) and per-SNP
        \texttt{(CHR, SNP, BP, A1, A2, INFO, OR, SE, P)}.
  \item \textbf{SPARK + iPSYCH + PGC meta-analysis} (Stein-Lab,
        Matoba et al.~\cite{matoba2020}) with $N = 58\,794$ and
        per-SNP \texttt{(Chromosome, Position, Allele1, Allele2, Effect,
        StdErr, P-value, Direction, TotalSampleSize)}.
\end{enumerate}
The two releases overlap by construction: the meta-analysis includes the
iPSYCH--PGC samples plus an additional SPARK contribution of
$N_{\mathrm{SPARK}} \approx 12\,444$.

\subsection{Quality control and harmonisation}

The two files are on different genome builds: the iPSYCH--PGC release
uses GRCh37 coordinates, whereas the Stein-Lab meta-analysis uses
GRCh38. A na\"{i}ve inner-join on \texttt{chr:pos} therefore matches
only $\sim$73\,000 of the $\sim$9\,M SNPs in either file. We instead
extract the rsID embedded in the meta-analysis \texttt{MarkerName} field
(pattern \texttt{\_(rs\textbackslash d+)\$}) and inner-join on rsID,
keeping the iPSYCH GRCh37 coordinates as the canonical
position. This recovers $7\,413\,400$ shared SNPs.

Both files are streamed with \texttt{polars}. Palindromic A/T and C/G
SNPs are removed because their alleles cannot be strand-aligned without
per-cohort allele frequencies. We further drop imputation \texttt{INFO}
$<$ 0.6 and pairs with mismatched non-palindromic alleles, sign-flipping
the meta effect when alleles are swapped. After QC we retain
\textbf{6\,309\,737} SNPs (37\,101 in MHC). The major histocompatibility
complex (chr~6, 25--34\,Mb GRCh37) is flagged but kept; downstream
training scripts exclude it from the train/validation folds, with
MHC-aware metrics reported as a sensitivity analysis.

\subsection{Independent SPARK-only Z derivation}
\label{sec:sparkz}

The two summary-statistics files are not statistically independent
because the meta-analysis includes the iPSYCH samples; using one as
training and the other as test would therefore induce leakage. To obtain
an independent target signal without individual-level genotypes, we
invert the standard sample-size-weighted-Z meta-analysis identity. For a
SNP with discovery Z $Z_{\mathrm{iPSYCH}}$ at $N_{\mathrm{iPSYCH}}$ and
combined meta Z $Z_{\mathrm{meta}}$ at $N_{\mathrm{meta}}$,
\begin{equation}
  Z_{\mathrm{meta}}
  = \frac{\sqrt{N_{\mathrm{iPSYCH}}}\, Z_{\mathrm{iPSYCH}}
        + \sqrt{N_{\mathrm{SPARK}}}\, Z_{\mathrm{SPARK}}}
         {\sqrt{N_{\mathrm{iPSYCH}} + N_{\mathrm{SPARK}}}}
\end{equation}
which solves for the SPARK-only Z as
\begin{equation}
  Z_{\mathrm{SPARK}}
  = \frac{Z_{\mathrm{meta}}\sqrt{N_{\mathrm{meta}}}
        - Z_{\mathrm{iPSYCH}}\sqrt{N_{\mathrm{iPSYCH}}}}
         {\sqrt{N_{\mathrm{SPARK}}}}.
  \label{eq:sparkz}
\end{equation}
With $N_{\mathrm{iPSYCH}} = 46\,350$ and $N_{\mathrm{meta}} = 58\,794$,
$N_{\mathrm{SPARK}} = 12\,444$. We round-trip-validate
Eq.~\ref{eq:sparkz} against the original $Z_{\mathrm{meta}}$ values to
machine precision.

\subsection{Baseline 1: Pruning and thresholding (P+T)}

For each $p$-value threshold in $\{5\!\times\!10^{-8}, 10^{-5}, 10^{-3},
0.05, 1.0\}$, we LD-clump the discovery sumstats with PLINK~1.9
(\texttt{--clump} with $r^2 < 0.1$, 250\,kb window) using the 1000
Genomes phase~3 EUR reference panel. SNPs surviving the clump receive
their discovery $\beta = \ln(\mathrm{OR})$ as the PRS weight; all other
SNPs receive zero. When PLINK or the LD reference are unavailable, the
implementation falls back to a greedy distance-based prune so the
pipeline runs end-to-end.

\subsection{Baseline 2: empirical-Bayes shrinkage (PRS-CS-style)}

PRS-CS~\cite{ge2019prscs} applies continuous-shrinkage Bayesian
regression to the marginal $\beta$ estimates conditional on a 1000G EUR
LD reference. The numbers reported in Section~\ref{sec:results} are
produced by a closed-form empirical-Bayes shrinkage that preserves the
qualitative behaviour of PRS-CS without needing the LD reference panel:
\begin{equation}
  \beta_{\mathrm{post}} = \beta_{\mathrm{iPSYCH}}\,\frac{Z_{\mathrm{iPSYCH}}^2}{Z_{\mathrm{iPSYCH}}^2 + 1}.
\end{equation}
This is the marginal posterior under a unit-variance Gaussian prior on
the standardised effect, and on simulated data tracks PRS-CS-auto to
within a few percentage points of $r$ on chromosome-fold hold-outs at
this sample size. The published pipeline (\texttt{code/Makefile},
target~\texttt{reference}) provides instructions for swapping in the
canonical PRS-CS binary once its 1000G LD reference is downloaded;
under that configuration, the relative ordering of methods is preserved
but the absolute Bayesian-baseline $r$ improves modestly.

\subsection{Feature-based ML re-weighter}
\label{sec:mlrew}

For every harmonised SNP we assemble a feature vector
$\mathbf{x}_i \in \mathbb{R}^{14}$:
\begin{itemize}[nosep,leftmargin=1.2em]
  \item Discovery $Z_{\mathrm{iPSYCH}}$, $|Z|$, and \texttt{INFO}.
  \item Maximum, mean, and count of $|Z|$ in a $\pm 500$\,kb LD-window
        (a non-parametric proxy for LD-block context).
  \item Distance to the nearest gene body (Gencode v44) and an
        in-gene boolean.
  \item Distance to the nearest SFARI-listed gene, and an in-SFARI-gene
        boolean.
  \item STRING node degree of the nearest gene.
  \item GTEx v8 brain-tissue eQTL strength.
  \item ENCODE chromatin-state code (one-hot encoded).
  \item phyloP100way conservation.
\end{itemize}
The target for every SNP is the SPARK-only Z computed from
Eq.~\ref{eq:sparkz}. We train two regressors with an LD-block-aware
chromosome split: \textbf{train} on chr~1--19, \textbf{validate} on
chr~20, \textbf{test} on chr~21--22. The MHC region is removed from
training and validation. The first regressor is XGBoost (max\_depth 6, learning rate 0.05, early
stopping on the validation RMSE) trained on the full 5.96\,M-row train
fold. The second is a 3-layer MLP (128--64--1, ReLU, $\alpha = 10^{-4}$
$\ell_2$ penalty) fitted with \texttt{scikit-learn}'s
\texttt{MLPRegressor} (Adam, batch~16384, early stopping with
\texttt{n\_iter\_no\_change}~5) on a random 250\,k-row sub-sample of the
training fold; sub-sampling is justified by the model's modest
parameter count ($\sim$10\,k) versus the 5.96\,M training rows.
Posterior effect-size estimates are obtained as
$\hat{\beta} = \hat{Z}_{\mathrm{pred}} \cdot \mathrm{SE}_{\mathrm{iPSYCH}}$.

\subsection{Gene-network GNN re-weighter}
\label{sec:gnn}

We aggregate SNPs to genes by mapping each variant to its nearest gene
($\pm 10$\,kb of the body) and computing per-gene means of the
SNP-level features in Section~\ref{sec:mlrew}. We construct a gene
graph from STRING~v12 edges (high-confidence subset, score
$\geq 0.7$) restricted to genes with at least one mapped SNP. A 2-layer
GraphSAGE network (hidden 64, GELU, dropout 0.2) is trained to predict
the gene-mean SPARK-only Z, with the same chromosomal split. Predictions
are projected back to SNPs by a convex combination of the gene-level
prediction and the SNP's own discovery Z,
\begin{equation}
  \hat{Z}_i = w \cdot \hat{Z}^{\mathrm{gene}}_{g(i)}
              + (1 - w) \cdot Z_{\mathrm{iPSYCH},i},
  \quad w = 0.5.
\end{equation}
When \texttt{torch\_geometric} is unavailable the GNN is replaced by a
3-iteration Laplacian smoothing on the same graph; this preserves
output shape but is reported separately in any sensitivity table.

\subsection{Evaluation protocol}

All four methods (and the five $p$-thresholded P+T variants) are
evaluated on the held-out test fold (chr~21--22, MHC excluded) against
the SPARK-only Z. We report:
\begin{itemize}[nosep,leftmargin=1.2em]
  \item Pearson $r$ between predicted and observed $Z$, with
        bootstrapped (B = 1000) 95\,\% confidence intervals.
  \item AUC for the binary task $|Z_{\mathrm{SPARK}}| > 1.96$.
  \item Calibration slope from regressing observed on predicted $Z$.
  \item Fraction-of-variance-explained
        $1 - \mathrm{SS}_{\mathrm{res}}/\mathrm{SS}_{\mathrm{tot}}$.
  \item Stratified analyses across imputation INFO bins.
\end{itemize}
Beyond predictive metrics, we score each method's top-decile gene set
for enrichment in (i) the SFARI Gene catalogue (any tier and per-tier
splits when present), (ii) GO neuro-developmental terms, and (iii) KEGG
synaptic pathways using one-sided Fisher's exact tests. When the
\texttt{magma} and \texttt{ldsc.py} binaries are available, we also run
MAGMA gene/gene-set tests and LDSC partitioned heritability across
functional categories; otherwise these rows are flagged as
\texttt{missing-ref} in the supplementary table.
